\appendix
\section{Secondary results}
\label{sec:appendix}

This appendix reports results that support the main text but carry none of its claims.

\subsection{Sensitivity to evidence orientation}
\label{sec:orientation}

The sign convention of~\eqref{eq:orient} deserves its own analysis, because
getting it wrong produces a detector that passes the tests most papers run and
fails the one they do not. The transform~\eqref{eq:phi} is even, but fusion
happens before $\phi$, so it is the summed statistic that is folded, not the
individual branches.

Suppose the two statistics were summed without orientation, for a beacon whose
behavioural anomaly reaches $|z_B| = 3$ and whose localisation anomaly reaches
$z_L = 3$ at full intensity, with equal weights. Under A1, $z_B \approx 0$ and
the misoriented fusion gives $Z = +2.12$, hence $T^D = 0.10$: confirmed. Under
A2, $z_L \approx 0$ and $Z = -2.12$, which because $\phi$ is even also gives
$T^D = 0.10$: confirmed. Both single-axis attacks are caught. Under a hybrid
attack at half intensity, however, $z_B \approx -1.5$ and $z_L \approx +1.5$,
and
\begin{equation}
\label{eq:cancel}
Z = \frac{(-1.5) + (+1.5)}{\sqrt2} = 0, \qquad T^D = \phi(0) = 1.0.
\end{equation}
The attacker's two misbehaviours cancel exactly, and a node misbehaving on both
axes receives perfect trust. A detector with this defect is blind to hybrid
attacks in proportion to how evenly the attacker balances the two axes, and an
adversary aware of the defect could balance them deliberately. With
\eqref{eq:orient} applied, the same attacker gives
$Z = \bigl((+1.5) + (+1.5)\bigr)/\sqrt2 = 2.12$, recovering the full-intensity
single-axis statistic from two half-intensity observations.


\subsection{Localisation error}
\label{sec:locerr}

Throughput is uninformative for A1, since the attack has no packet-level
footprint. The service A1 degrades is localisation, so we report the
displacement between the position a beacon advertises and its true position,
which is the error a consumer of the localisation service inherits
(Table~\ref{tab:locerr}).

\begin{table}[!t]
\caption{Localisation error in metres, means across five seeds of the original
campaign. ``Attacker'' covers receptions from the compromised beacon; ``all''
covers every monitored beacon, of which two of three are honest.}
\label{tab:locerr}
\centering
\small
\begin{tabular}{@{}llrrrr@{}}
\toprule
Scenario & Scope & RMSE & MAE & Median & P95 \\
\midrule
Clean & all & 0.05 & 0.02 & 0.00 & 0.06 \\
A1 & attacker & 267.25 & 265.42 & 271.41 & 299.23 \\
   & all & 145.27 & 78.78 & 0.01 & 288.77 \\
A2 & attacker & 0.03 & 0.01 & 0.00 & 0.05 \\
   & all & 0.04 & 0.01 & 0.00 & 0.05 \\
A3 & attacker & 146.02 & 145.81 & 145.91 & 150.01 \\
   & all & 77.88 & 41.66 & 0.01 & 150.00 \\
\bottomrule
\end{tabular}
\end{table}

The clean-network figure of 0.05\,m RMSE confirms that honest beacons advertise
their true positions to within numerical precision, so any error observed
elsewhere is attributable to the attack. A2 produces 0.04\,m, indistinguishable
from clean: A2 is a pure forwarding attack with no positional signature, exactly
as A1 is a pure positional attack with no forwarding signature. A3's attacker
error of 146.02\,m against A1's 267.25\,m confirms the half-intensity
construction. The ratio is 0.55 instead of 0.50 because clamping binds less
often at the smaller offset.

The substantive result is the ``all'' row. Under A1 the mean error across all
monitored beacons is 145.27\,m while the median is 0.01\,m. Two of three beacons
are honest and contribute essentially zero error, while the third contributes
the full displacement. A localisation consumer drawing on the affected beacon
inherits an error more than three orders of magnitude above the clean-network floor, and
averaging over honest neighbours does not conceal it, since trilateration
depends on all contributing anchors.


\subsection{Mechanism cost}

\begin{figure}[!t]
\centering
\includegraphics[width=0.8\textwidth]{fig7_mechanism}
\caption{Isolation mechanism accounting in one illustrative enforced run (A2,
seed 1), not an aggregate. Policy D suppresses more attacker drops but blocks
more than twice as many frames and performs three times as many isolation
events as policy E.}
\label{fig:mechanism}
\end{figure}

Fig.~\ref{fig:mechanism} decomposes what each policy does within a single run.
The counters were logged per run but not aggregated across the campaign, so the
figure illustrates the accounting without establishing a general magnitude.
Both policies suppress attacker drops substantially relative to the control:
from 750 packets under C to 69 under D and 111 under E. Frames blocked by the
enforcement mechanism itself differ by more than a factor of two, 3381 under D
against 1566 under E, tracking isolation-event counts of 72 and 24. The counter
separates the mechanism's benefit from its cost within a run, which aggregate
PDR does not.


\subsection{Aggregation ablation}
\label{sec:ablation}

\begin{table}[!t]
\caption{Aggregation ablation, original campaign. Confirmations out of 15
target-seeds per cell.}
\label{tab:ablation}
\centering
\small
\begin{tabular}{@{}lcccc@{}}
\toprule
Configuration & A1 & A2 & A3 & Clean \\
\midrule
Direct only & 6/15 & 5/15 & 6/15 & 1/15 \\
Direct + indirect & 6/15 & 5/15 & 6/15 & 1/15 \\
Direct + indirect + historical & 6/15 & 5/15 & 6/15 & 1/15 \\
\bottomrule
\end{tabular}
\par\smallskip\footnotesize
Counts include false confirmations: A1 6/15 is the attacker in five seeds plus
one honest beacon.
\end{table}

\begin{figure}[!t]
\centering
\includegraphics[width=0.62\textwidth]{fig9_ablation}
\caption{Aggregation ablation. The three configurations are identical in every
cell: adding indirect and historical trust changes no verdict.}
\label{fig:ablation}
\end{figure}

Table~\ref{tab:ablation} and Fig.~\ref{fig:ablation} report the second stage
ablated across its components. History-only is not an admissible configuration,
since $T^H$ seeds from the combined value. The three configurations are
identical in every cell. The components are not inert in the underlying
computation: on A2 seed~1, direct-only aggregation gives mean trust 0.7148
with 50 window-level confirmations, while direct plus indirect gives 0.7548
with 70. But the differences never cross the persistence gate, so no target's
final verdict changes. $T^H$ is an exponentially weighted average of direct
trust and $T^I$ a trimmed mean of peers' direct trust: both are smoothed
functions of the same underlying signal, and they would be expected to diverge
only where direct trust is noisy and observers disagree, a regime our topology
does not produce often enough to matter.


